\documentclass[a4paper,12pt]{article}

\usepackage[T2A]{fontenc}
\usepackage[utf8]{inputenc}
\usepackage[russian]{babel}

\usepackage{amsmath,amssymb,amsthm,mathtools}
\usepackage{helvet}
\renewcommand{\familydefault}{\sfdefault}
\usepackage{icomma}

\usepackage[a4paper,margin=2.1cm,headheight=15pt]{geometry}
\usepackage{microtype}

\usepackage{tikz}
\usepackage{tkz-euclide}
\usepackage{pgfplots}
\pgfplotsset{compat=1.18}

\usepackage{tabularx,booktabs,longtable}
\usepackage{enumitem}
\usepackage{hyperref}
\usepackage{parskip}
\usepackage{fancyhdr}
\usepackage{setspace}
\usepackage{xcolor}

% -------------------------------------------------
% Цветовая палитра
% -------------------------------------------------
\definecolor{accent}{HTML}{008080}
\definecolor{darkaccent}{HTML}{004D4D}
\definecolor{textgray}{HTML}{333333}
\definecolor{softgray}{HTML}{EEF3F3}
\definecolor{linegray}{HTML}{B8C7C7}

% -------------------------------------------------
% Общая типографика
% -------------------------------------------------
\onehalfspacing
\setlength{\parindent}{0pt}
\setlength{\parskip}{0.65em}

\hypersetup{
    colorlinks=true,
    linkcolor=darkaccent,
    urlcolor=darkaccent
}

\pagestyle{fancy}
\fancyhf{}
\fancyhead[L]{\small\textcolor{textgray}{Текстовые задачи на движение}}
\fancyhead[R]{\small\textcolor{textgray}{Кирилл Зиновьев}}
\fancyfoot[C]{\small\textcolor{textgray}{\thepage}}
\renewcommand{\headrulewidth}{0.3pt}
\renewcommand{\headrule}{
    \hbox to\headwidth{\color{linegray}\leaders\hrule height \headrulewidth\hfill}
}

\setlist[itemize]{
    leftmargin=1.5em,
    itemsep=0.25em,
    topsep=0.3em
}

\setlist[enumerate]{
    leftmargin=2em,
    itemsep=0.55em,
    topsep=0.4em
}

\newcommand{\km}{\text{ км}}
\newcommand{\kmh}{\text{ км/ч}}
\newcommand{\hour}{\text{ ч}}
\newcommand{\minutes}{\text{ мин}}

\newcommand{\sectionline}{
    \vspace{-0.2em}
    {\color{linegray}\rule{\linewidth}{0.5pt}}
    \vspace{0.5em}
}

\newcommand{\key}[1]{\textcolor{darkaccent}{\textbf{#1}}}

% =================================================
\begin{document}

% =================================================
% ТИТУЛЬНЫЙ ЛИСТ
% =================================================
\begin{titlepage}
\thispagestyle{empty}

\begin{center}

\vspace*{2.6cm}

{\Large\textcolor{darkaccent}{\textbf{ЧЕК-ЛИСТ}}}

\vspace{1.0cm}

{\Huge\bfseries
Текстовые задачи\\[0.2em]
на движение
}

\vspace{0.6cm}

{\Large\textcolor{accent}{
Таблица \(v\)--\(t\)--\(S\), схема движения\\
и относительная скорость
}}

\vspace{2.2cm}

{\large
Лёвин Артём Александрович\\[0.25em]
\textcolor{darkaccent}{\textbf{эксклюзивно для Кирилла Зиновьева}}
}

\vfill

{\large \today}

\vspace{2.4cm}

\end{center}

% Сдержанная кривая Безье внизу титульного листа
\begin{tikzpicture}[remember picture,overlay]
    \draw[
        color=accent,
        line width=1.1pt
    ]
    ([xshift=1.9cm,yshift=1.55cm]current page.south west)
    .. controls
    ([xshift=6.0cm,yshift=2.35cm]current page.south west)
    and
    ([xshift=-6.0cm,yshift=0.85cm]current page.south east)
    ..
    ([xshift=-1.9cm,yshift=1.55cm]current page.south east);
\end{tikzpicture}

\end{titlepage}

% =================================================
% ОСНОВНОЙ ЧЕК-ЛИСТ
% =================================================

\section*{\textcolor{darkaccent}{1. Главная идея}}

\sectionline

Большинство задач на движение удобно переводить из текста в
\key{математическую модель} в три этапа:

\[
\boxed{\text{таблица}}
\quad\longrightarrow\quad
\boxed{\text{схема}}
\quad\longrightarrow\quad
\boxed{\text{уравнение}}.
\]

Таблица помогает распределить данные, схема показывает взаимное
расположение объектов, а уравнение связывает известные и неизвестные
величины.

Главные величины:

\[
\boxed{S=vt},
\qquad
\boxed{v=\frac{S}{t}},
\qquad
\boxed{t=\frac{S}{v}},
\]

где

\[
S\text{ --- путь},\qquad
v\text{ --- скорость},\qquad
t\text{ --- время}.
\]

\textbf{Важно.}
Перед вычислениями единицы измерения должны быть согласованы.
Например, при скорости в \(\kmh\) время удобно выражать в часах.

% -------------------------------------------------

\section*{\textcolor{darkaccent}{2. Таблица \(v\)--\(t\)--\(S\)}}

\sectionline

Для двух участников движения удобно использовать таблицу:

\begin{table}[h!]
\centering
\caption{Базовая таблица движения}
\renewcommand{\arraystretch}{1.35}
\begin{tabularx}{0.88\linewidth}{Xccc}
\toprule
& \textbf{Скорость \(v\)} & \textbf{Время \(t\)} & \textbf{Путь \(S\)}\\
\midrule
Первый объект  & \(v_1\) & \(t_1\) & \(v_1t_1\)\\
Второй объект  & \(v_2\) & \(t_2\) & \(v_2t_2\)\\
\bottomrule
\end{tabularx}
\end{table}

Если участники начали движение \key{одновременно} и рассматривается один
и тот же момент времени, то

\[
t_1=t_2.
\]

Если неизвестное время обозначено через \(x\), получаем:

\[
S_1=v_1x,
\qquad
S_2=v_2x.
\]

\textbf{Пример.}
Если скорости пешеходов равны \(4\kmh\) и \(5\kmh\), а каждый идёт
\(3\hour\), то

\[
S_1=4\cdot3=12\km,
\qquad
S_2=5\cdot3=15\km.
\]

% -------------------------------------------------

\section*{\textcolor{darkaccent}{3. Встречное движение}}

\sectionline

Пусть два объекта находятся на расстоянии \(S_0\) друг от друга и
движутся \key{навстречу}.

\begin{center}
\begin{tikzpicture}[x=0.12cm,y=1cm,>=stealth]
    \draw[line width=0.8pt] (0,0) -- (70,0);

    \fill[accent] (0,0) circle (2pt);
    \fill[accent] (70,0) circle (2pt);

    \node[below=5pt] at (0,0) {первый};
    \node[below=5pt] at (70,0) {второй};

    \draw[->,thick,accent] (5,0.35) -- (25,0.35);
    \draw[->,thick,accent] (65,0.35) -- (45,0.35);

    \node[above] at (15,0.35) {\(v_1\)};
    \node[above] at (55,0.35) {\(v_2\)};

    \draw[<->,line width=0.6pt] (0,-0.65) -- (70,-0.65);
    \node[fill=white] at (35,-0.65) {\(S_0\)};
\end{tikzpicture}
\end{center}

За время \(t\) они пройдут:

\[
S_1=v_1t,
\qquad
S_2=v_2t.
\]

Расстояние между ними уменьшится на

\[
S_1+S_2.
\]

Поэтому оставшееся расстояние:

\[
\boxed{S_{\text{ост}}=S_0-(v_1+v_2)t}.
\]

Если объекты встретились, то \(S_{\text{ост}}=0\), следовательно,

\[
\boxed{t_{\text{встр}}=\frac{S_0}{v_1+v_2}}.
\]

% -------------------------------------------------

\section*{\textcolor{darkaccent}{4. Скорость сближения}}

\sectionline

\key{Скорость сближения} показывает, на сколько единиц расстояния между
объектами уменьшается за единицу времени.

Для движения навстречу:

\[
\boxed{v_{\text{сбл}}=v_1+v_2}.
\]

Например,

\[
v_1=4\kmh,\qquad v_2=5\kmh,
\]

поэтому

\[
v_{\text{сбл}}=4+5=9\kmh.
\]

За \(3\hour\) расстояние уменьшится на

\[
9\cdot3=27\km.
\]

Если первоначально между пешеходами было \(40\km\), то останется

\[
40-27=13\km.
\]

\textbf{Это тот же результат}, который получается через таблицу:

\[
12+13+15=40.
\]

% -------------------------------------------------

\section*{\textcolor{darkaccent}{5. Если нужно найти время}}

\sectionline

Пусть первоначальное расстояние равно \(40\km\), скорости ---
\(4\kmh\) и \(5\kmh\), а требуется узнать, через какое время между
пешеходами останется \(13\km\).

Обозначим время через \(x\).

\begin{table}[h!]
\centering
\caption{Перевод условия в математическую модель}
\renewcommand{\arraystretch}{1.35}
\begin{tabularx}{0.88\linewidth}{Xccc}
\toprule
& \(v\) & \(t\) & \(S\)\\
\midrule
Первый пешеход & \(4\) & \(x\) & \(4x\)\\
Второй пешеход & \(5\) & \(x\) & \(5x\)\\
\bottomrule
\end{tabularx}
\end{table}

Получаем:

\[
4x+13+5x=40.
\]

Приводим подобные слагаемые:

\[
9x+13=40,
\]

\[
9x=27,
\]

\[
\boxed{x=3\hour}.
\]

Здесь используется важное алгебраическое правило:

\[
4x+5x=9x,
\]

поскольку слагаемые имеют одинаковую буквенную часть.

Но, например,

\[
25+3x
\]

сложить в одно слагаемое нельзя.

% -------------------------------------------------

\section*{\textcolor{darkaccent}{6. Движение в одном направлении}}

\sectionline

Если оба объекта движутся в одну сторону и более быстрый объект
догоняет более медленный, расстояние между ними уменьшается со
скоростью

\[
\boxed{v_{\text{сбл}}=v_{\text{быстр}}-v_{\text{медл}}}.
\]

\begin{center}
\begin{tikzpicture}[x=0.13cm,y=1cm,>=stealth]
    \draw[line width=0.8pt] (0,0) -- (65,0);

    \fill[accent] (10,0) circle (2pt);
    \fill[accent] (45,0) circle (2pt);

    \draw[->,thick,accent] (12,0.35) -- (30,0.35);
    \draw[->,thick,accent] (47,0.35) -- (58,0.35);

    \node[above] at (21,0.35) {\(v_{\text{быстр}}\)};
    \node[above] at (52.5,0.35) {\(v_{\text{медл}}\)};

    \draw[<->,line width=0.6pt] (10,-0.65) -- (45,-0.65);
    \node[fill=white] at (27.5,-0.65) {\(S_0\)};
\end{tikzpicture}
\end{center}

Если начальное расстояние между ними \(S_0\), то время догоняния:

\[
\boxed{
t_{\text{дог}}=
\frac{S_0}{v_{\text{быстр}}-v_{\text{медл}}}
}.
\]

Формула имеет смысл только при

\[
v_{\text{быстр}}>v_{\text{медл}}.
\]

% -------------------------------------------------

\section*{\textcolor{darkaccent}{7. Сближение и удаление --- не одно и то же}}

\sectionline

Нельзя пользоваться правилом
«разные стороны --- складываем» без анализа схемы.

Если два объекта движутся \key{навстречу}, расстояние между ними
уменьшается:

\[
S(t)=S_0-(v_1+v_2)t.
\]

Если они движутся \key{в противоположные стороны друг от друга},
расстояние между ними увеличивается:

\[
\boxed{S(t)=S_0+(v_1+v_2)t}.
\]

В обоих случаях модуль относительной скорости равен

\[
v_1+v_2,
\]

но в первом случае происходит \key{сближение}, а во втором ---
\key{удаление}.

Именно поэтому перед выбором действия сначала полезно нарисовать
схему со стрелками.

% -------------------------------------------------

\section*{\textcolor{darkaccent}{8. Универсальный алгоритм}}

\sectionline

\begin{enumerate}[label=\textbf{\arabic*.}]
    \item Внимательно определите, \key{кто}, \key{откуда} и
    \key{в каком направлении} движется.

    \item Выпишите известные величины и их единицы измерения.

    \item Если возможно, составьте таблицу
    \[
    v,\quad t,\quad S.
    \]

    \item Неизвестную величину обозначьте буквой, обычно \(x\).

    \item Недостающие пути выразите по формуле
    \[
    S=vt.
    \]

    \item Нарисуйте схему и отметьте на ней:
    \begin{itemize}
        \item начальное расстояние;
        \item пройденные участки;
        \item оставшееся расстояние;
        \item направления движения.
    \end{itemize}

    \item По схеме составьте уравнение.

    \item Решите уравнение и проверьте:
    \begin{itemize}
        \item знак ответа;
        \item размерность;
        \item соответствие условию.
    \end{itemize}

    \item Запишите ответ с единицами измерения.
\end{enumerate}

% -------------------------------------------------

\section*{\textcolor{darkaccent}{9. Два способа решения}}

\sectionline

Одну и ту же задачу часто можно решить двумя способами.

\begin{table}[h!]
\centering
\caption{Сравнение подходов}
\renewcommand{\arraystretch}{1.35}
\begin{tabularx}{\linewidth}{p{0.24\linewidth}XX}
\toprule
\textbf{Метод} & \textbf{Преимущество} & \textbf{Когда особенно полезен}\\
\midrule
Таблица + схема + уравнение
&
Все данные видны; меньше риск потерять условие; формируется универсальный способ решения
&
При новых, сложных и экзаменационных задачах
\\
\addlinespace
Относительная скорость
&
Позволяет получить короткое решение
&
Для простых задач, проверки результата и быстрого счёта
\\
\bottomrule
\end{tabularx}
\end{table}

Для уверенного решения задач полезно владеть \key{обоими способами}.

% -------------------------------------------------

\section*{\textcolor{darkaccent}{10. Типичные ошибки}}

\sectionline

\begin{itemize}
    \item \textbf{Сложить скорости автоматически.}

    Сначала определите, уменьшается или увеличивается расстояние между
    объектами.

    \item \textbf{Вычесть скорости при встречном движении.}

    При движении навстречу:
    \[
    v_{\text{сбл}}=v_1+v_2.
    \]

    \item \textbf{Считать время одинаковым, если участники начали
    движение не одновременно.}

    При задержке одного участника времена движения различаются.

    \item \textbf{Забыть оставшееся расстояние.}

    Если встреча ещё не произошла, обычно
    \[
    S_1+S_{\text{ост}}+S_2=S_0.
    \]

    \item \textbf{Складывать неподобные слагаемые.}

    Верно:
    \[
    4x+5x=9x.
    \]

    Неверно:
    \[
    25+3x=28x.
    \]

    \item \textbf{Потерять единицы измерения.}

    Ответ \(5\) сам по себе неполон. Нужно записать, например,
    \(5\kmh\) или \(5\hour\).

    \item \textbf{Не проверить смысл результата.}

    В обычной задаче на движение время, скорость и путь не могут
    оказаться отрицательными.
\end{itemize}

% =================================================
% ТРЕНИРОВОЧНЫЙ БЛОК
% =================================================

\clearpage

\section*{\textcolor{darkaccent}{Тренировочный блок}}

\sectionline

Для каждой задачи рекомендуется:

\[
\boxed{\text{таблица}}
\longrightarrow
\boxed{\text{схема}}
\longrightarrow
\boxed{\text{уравнение}}
\longrightarrow
\boxed{\text{ответ}}.
\]

\subsection*{\textcolor{accent}{Средний уровень}}

\begin{enumerate}[label=\textbf{\arabic*.}]

\item
Из двух посёлков, расстояние между которыми \(42\km\), одновременно
вышли навстречу друг другу два пешехода со скоростями
\(5\kmh\) и \(7\kmh\).
Через сколько часов они встретятся?

\item
Из двух пунктов, расстояние между которыми \(50\km\), одновременно
навстречу друг другу вышли два пешехода со скоростями
\(4\kmh\) и \(6\kmh\).
Какое расстояние будет между ними через \(3\hour\)?

\item
Два велосипедиста движутся в одном направлении.
Первый едет впереди со скоростью \(5\kmh\), второй ---
со скоростью \(8\kmh\).
Первоначальное расстояние между ними равно \(12\km\).
Через сколько часов второй велосипедист догонит первого?

\end{enumerate}

\subsection*{\textcolor{accent}{Выше среднего}}

\begin{enumerate}[label=\textbf{\arabic*.},start=4]

\item
Из двух городов, расстояние между которыми \(68\km\), одновременно
навстречу друг другу вышли два туриста.
Скорость первого равна \(7\kmh\).
Через \(4\hour\) расстояние между ними стало равно \(16\km\).
Найдите скорость второго туриста.

\item
Расстояние между двумя населёнными пунктами равно \(90\km\).
Из них одновременно навстречу друг другу выехали два велосипедиста со
скоростями \(8\kmh\) и \(6\kmh\).
Через сколько часов расстояние между ними станет равно \(27\km\)?

\item
Из пункта \(A\) выехал велосипедист со скоростью \(12\kmh\).
Через \(1\hour\) вслед за ним выехал второй велосипедист со скоростью
\(16\kmh\).
Через сколько часов после своего выезда второй велосипедист догонит
первого?
На каком расстоянии от пункта \(A\) произойдёт встреча?

\item
Два пешехода одновременно вышли из одного места в противоположных
направлениях со скоростями \(4{,}5\kmh\) и \(5{,}5\kmh\).
Через сколько часов расстояние между ними станет равно \(25\km\)?

\item
Два туриста находятся на расстоянии \(75\km\) друг от друга и начинают
одновременно двигаться в противоположных направлениях друг от друга.
Их скорости равны \(6\kmh\) и \(4\kmh\).
Какое расстояние будет между ними через \(2\hour\)?

\item
Два автомобиля движутся в одном направлении.
Впереди едет автомобиль со скоростью \(6\kmh\).
Второй находится на \(18\km\) позади него и догоняет первый за
\(3\hour\).
Найдите скорость второго автомобиля.

\end{enumerate}

\subsection*{\textcolor{accent}{Сложная задача}}

\begin{enumerate}[label=\textbf{\arabic*.},start=10]

\item
Расстояние между двумя посёлками равно \(126\km\).
Из первого посёлка в сторону второго вышел пешеход со скоростью
\(6\kmh\).
Через \(1{,}5\hour\) из второго посёлка навстречу ему вышел другой
пешеход со скоростью \(8\kmh\).

Через сколько часов после выхода второго пешехода расстояние между ними
станет равно \(33\km\)?

\end{enumerate}

% =================================================
% ОТВЕТЫ
% =================================================

\clearpage

\section*{\textcolor{darkaccent}{Ответы к тренировочному блоку}}

\sectionline

\begin{table}[h!]
\centering
\caption{Ответы}
\renewcommand{\arraystretch}{1.45}
\begin{tabularx}{0.92\linewidth}{cX}
\toprule
\textbf{№} & \textbf{Ответ}\\
\midrule

1 &
\(3{,}5\hour\)
\\

2 &
\(20\km\)
\\

3 &
\(4\hour\)
\\

4 &
\(6\kmh\)
\\

5 &
\(4{,}5\hour\)
\\

6 &
\(3\hour\) после выезда второго велосипедиста;
встреча произойдёт в \(48\km\) от пункта \(A\)
\\

7 &
\(2{,}5\hour\)
\\

8 &
\(95\km\)
\\

9 &
\(12\kmh\)
\\

10 &
\(6\hour\) после выхода второго пешехода
\\

\bottomrule
\end{tabularx}
\end{table}

\vfill

\begin{center}
\textcolor{darkaccent}{
\textbf{Главная памятка}
}

\vspace{0.5em}

\[
\boxed{S=vt}
\]

\[
\boxed{
\text{навстречу: }\;
v_{\text{сбл}}=v_1+v_2
}
\]

\[
\boxed{
\text{догонка: }\;
v_{\text{сбл}}=v_{\text{быстр}}-v_{\text{медл}}
}
\]

\vspace{0.8em}

\textit{
Сначала определите направления движения,
затем составляйте формулы.
}

\end{center}

\end{document}